\section{Zadanie 12}
\autor{Aleksander Wieczorek}
Czas obsługi i-tego klienta dany jest zmienn¡ losową o wartości oczekiwanej wynoszącej 2 minuty oraz wariancji równej 1 minucie. Zmienne opisujące każdego z
klientów są niezależne. Niech Y oznacza sumaryczny czas obsługi 50 klientów. Wyznacz $P(90 \le Y \le 110)$.
\subsection*{Rozwiązanie}
W tym zadaniu trzeba skorzystać z CTG. Ogranicza się to do tego, że po prostu musimy tak zmodyfikować 
daną nierówność, żeby wyszedł nam wzór z CTG. Wtedy aplikujemy twierdzenie i mamy wynik.\\
Policzmy potrzebne rzeczy:\\
$$Y_i - \text{czas obsługi i-tego klienta}$$
$$Y=\sum_{i=1}^{50} Y_i \qquad \bar{Y} = \frac{Y}{50}$$
$$E(\bar{Y})=\frac{1}{50}E(Y)=\frac{1}{50} \cdot \sum_{i=1}^{50}E(Y_i)=\frac{1}{50} \cdot 50 \cdot 2=2$$
$$Var(\bar{Y})=\frac{1}{50^2}\cdot \sum_{i=1}^{50} Var(Y_i) =\frac{1}{50^2}\cdot 50 \cdot 1 = \frac{1}{50}$$
$$\frac{1}{\sqrt{Var(\bar Y)}}= \sqrt{50}=5\sqrt{2}$$
I teraz wystarczy odpowiednio przekształcić nierówność z treści:
$$P(90 \le Y \le 110) = P\left(\frac{9}{5}\le \bar{Y} \le \frac{11}{5}\right)= P\left(-\frac{1}{5}\le \bar{Y}-E(\bar{Y})\le \frac{1}{5} \right)$$
$$= P\left(-\sqrt{2} \le \frac{\bar{Y}-E(\bar{Y})}{\sqrt{Var(\bar Y )}}\le \sqrt{2} \right)=\Phi(\sqrt{2}) - \Phi(-\sqrt{2})=2\cdot\Phi(\sqrt{2})-1$$
Przedostatnia równość to CTG, a ostatnia wynika z symetrii rozkładu normalnego ($\Phi(-x)=1-\Phi(x)$).\\
Teraz możemy jedynie odczytać dystrybuantę z tabelki.
$$2\cdot\Phi(\sqrt{2})-1 \approx 2\cdot \Phi (1.41) - 1 \approx 0.84 $$

\subsection*{Odpowiedź}
\[
\boxed{\Phi(\sqrt{2}) - \Phi(-\sqrt{2}) \approx 0.84 }
\]
